\subsection{A matrix-factorization map and a cancellation on ramified nodes}

\begin{lemma}[Proved: an additive matrix-factorization map to the Swan kernel]
\label{C:MF-Swan-transfer}
Let $S_0$ be a regular noetherian local ring, let
$g\in S_0$ be a nonzero nonunit, and put $R=S_0/(g)$.
Let $K_0(\operatorname{MCM}R)$ be the Grothendieck group
of finitely generated maximal Cohen--Macaulay $R$-modules,
with relations from all short exact sequences of such modules.
There is a homomorphism
\[
 \Delta_g:K_0(\operatorname{MCM}R)/\mathbb Z[R]
   \longrightarrow
   \ker\bigl(R_R(C_2)\longrightarrow R_{R[1/2]}(C_2)\bigr).
                                                               \tag{C.1.MF1}
\]
For a presentation over $S_0$
\[
 0\longrightarrow S_0^r\xrightarrow{\Phi}S_0^r
       \longrightarrow M\longrightarrow0,
 \qquad \Phi\Psi=\Psi\Phi=gI_r,
\]
its value is
\[
 \Delta_g(M)=
 \left[\begin{pmatrix}I_r&\bar\Phi\\
                        \bar\Psi&-I_r\end{pmatrix}\right]
       -r([\mathbf1]+[\mathrm{sgn}]),                 \tag{C.1.MF2}
\]
where bars denote reduction modulo $g$.
The displayed involution has acyclic two-periodic Tate
complex over $R$.
No injectivity or retraction assertion is part of this lemma.
\end{lemma}
\begin{proof}
A maximal Cohen--Macaulay $R$-module has projective
dimension one over $S_0$, unless it is zero.
Indeed its depth over $S_0$ is $\dim S_0-1$, and
the Auslander--Buchsbaum formula over the regular local
ring applies.  This is the length-one presentation
argument in \cite[proof of Theorem 6.1, p.\ 53,
and Corollary 6.3, p.\ 55]{C-Eisenbud1980}; its hypotheses
do not require a coefficient field.
Since $gM=0$, multiplication by $g$ on the presenting
target factors uniquely through the injective $\Phi$.
The resulting map $\Psi$ satisfies $\Phi\Psi=gI$.
Multiplication on the right by $\Phi$ and cancellation
of $\Phi$ give $\Psi\Phi=gI$.
After inverting $g$, the source and target of $\Phi$
are isomorphic, so their finite free ranks agree.

Put
\[
 \sigma=\begin{pmatrix}I_r&0\\0&-I_r\end{pmatrix},
 \qquad D=\begin{pmatrix}0&\Phi\\\Psi&0\end{pmatrix}.
\]
Then $D^2=gI$, $\sigma D+D\sigma=0$, and
$(\sigma+D)^2=I+gI$ over $S_0$.
Thus its reduction $T$ is an involution on a free
$R$-module, as required.

The construction is independent of the presentation.
Changing its free bases changes $(\Phi,\Psi)$ by a
simultaneous matrix-factorization equivalence and conjugates
$T$ by the corresponding block-diagonal change of basis.
Any finite free presentation of length one is the sum
of a minimal presentation and finitely many identity
presentations: eliminate a unit matrix entry repeatedly.
Minimal finite free resolutions of isomorphic modules over
a local ring are isomorphic, by the comparison theorem
and Nakayama's lemma.
An identity summand of $\Phi$ has $g$ as the corresponding
summand of $\Psi$.  Its involution is
$\left(\begin{smallmatrix}1&1\\0&-1\end{smallmatrix}\right)$
over $R$ and is a regular $C_2$-lattice.
For example the second standard vector and its image
under the involution are an $R$-basis interchanged by it.
Its class is $[\mathbf1]+[\mathrm{sgn}]$, by the
augmentation sequence with a free sign kernel.
The subtraction in (C.1.MF2) cancels exactly this change.

For an exact sequence $0\to M'\to M\to M''\to0$,
apply the horseshoe construction to length-one finite
free $S_0$-resolutions of the endpoints.
The middle presentation has upper block matrix
\[
 \Phi_M=\begin{pmatrix}\Phi_{M'}&B\\0&\Phi_{M''}\end{pmatrix}.
\]
Its unique factorization partner $\Psi_M$ is also upper
block triangular with diagonal blocks $\Psi_{M'}$ and
$\Psi_{M''}$.
For example its lower-left block is zero because its
product with the injective $\Phi_{M''}$ is zero;
the remaining diagonal identities determine those blocks.
After regrouping the two endpoint summands, the middle
involution is upper block triangular with the endpoint
involutions on its diagonal.
It gives an admissible exact sequence of $R$-free
$C_2$-lattices.
The presentation ranks add, proving additivity of
(C.1.MF2).
For $M=R$, use the factorization $(g,1)$.
Its involution is again regular, so $\Delta_g(R)=0$.
This proves descent through the quotient in (C.1.MF1).

Over $R[1/2]$, $\bar D^2=0$ and
\[
 P=I+\bar D\sigma/2,
 \qquad P^{-1}=I-\bar D\sigma/2,
 \qquad P\sigma P^{-1}=\sigma+\bar D.
\]
Thus the constructed representation becomes exactly
$r$ free trivial lines and $r$ free sign lines.
This proves the kernel assertion without a rank-only
argument over $R[1/2]$.

Finally the two lifts $\widetilde T-I$ and
$\widetilde T+I$ have product $gI$ in both orders.
They are injective over $S_0$.
If, for example, $(\widetilde T-I)v=gw$, multiplying
by $\widetilde T+I$ and cancelling $g$ gives
$v=(\widetilde T+I)w$.
Reduction modulo $g$ proves exactness at one parity
of the Tate complex; the other is identical.
\end{proof}

\begin{proposition}[Proved: cancellation of the standard matrix-factorization lattice on every ramified exponent-three dyadic node]
\label{C:ramified-A2-MF-cancellation}
Let $A$ be a dyadic DVR with uniformizer $\pi$ and
$2=v\pi^e$, where $v\in A^\times$ and $e\geq2$.
Let $u\in A^\times$, and take either
\[
 R=\bigl(A[x,y]/(xy-u\pi^3)\bigr)_{(\pi,x,y)}
 \quad\text{or}\quad
 R=A[[x,y]]/(xy-u\pi^3).
\]
Write $z=y/u$, and put
\[
 \Phi=\begin{pmatrix}-x&-\pi^2\\\pi&z\end{pmatrix},
 \qquad
 \Psi=\begin{pmatrix}z&\pi^2\\-\pi&-x\end{pmatrix},
 \qquad
 T=\begin{pmatrix}I_2&\Phi\\\Psi&-I_2\end{pmatrix}.
                                                               \tag{C.1.MF3}
\]
Then $T$ is an involution and
\[
                    [T]=2([\mathbf1]+[\mathrm{sgn}])
                    \quad\text{in }R_R(C_2).          \tag{C.1.MF4}
\]
Its Tate complex is acyclic, but its underlying
$R[C_2]$-module is not projective.
This proves vanishing of this specified virtual class;
it does not prove surjectivity or failure of surjectivity
for the full node comparison map.
\end{proposition}
\begin{proof}
Put $g=\pi^3-xz$ in the corresponding regular local
ring $S_0$ before taking the quotient.
Direct multiplication gives $\Phi\Psi=\Psi\Phi=gI_2$,
so the preceding lemma applies.

The row $(\pi,x)$ presents $I_1=(\pi,x)$ as
$\operatorname{coker}\Phi$.
Here are the kernel details.
If $\pi a+xb=0$ in $R$, lift to $S_0$ and write
$\pi a+xb=gc$.  Then
\[
                 \pi(a-\pi^2c)+x(b+zc)=0.
\]
Since $(\pi,x)$ is a regular sequence in $S_0$,
there is $d\in S_0$ with
$(a-\pi^2c,b+zc)=(-xd,\pi d)$.
Thus $(a,b)=\Phi(d,-c)$ modulo $g$, as required.
Similarly the row $(x,\pi^2)$ presents
$I_2=(x,\pi^2)$ as $\operatorname{coker}\Psi$:
if $xa+\pi^2b=gc$, then
$x(a+zc)+\pi^2(b-\pi c)=0$, and the syzygies
of the regular sequence $(x,\pi^2)$ give
$(a,b)=\Psi(-c,-d)$ for a suitable $d$.
Both are maximal Cohen--Macaulay, either by these
matrix factorizations or their injective length-one
$S_0$-free presentations and the depth lemma.

We record two exact sequences of these actual ideals.
The first is the periodic-presentation sequence
\[
                    0\longrightarrow I_2\longrightarrow R^2
                       \xrightarrow{(\pi,x)}I_1\longrightarrow0.
                                                               \tag{C.1.MF5}
\]
Indeed its kernel is $\operatorname{im}\bar\Phi$,
which is $\operatorname{coker}\bar\Psi$ by
matrix-factorization exactness.
The second is
\[
 0\longrightarrow I_1\xrightarrow{c\mapsto(c,-\pi c)}
 R\oplus I_2\xrightarrow{(a,b)\mapsto\pi a+b}
 I_1\longrightarrow0.                              \tag{C.1.MF6}
\]
Surjectivity follows from $\pi R+I_2=I_1$.
For the kernel one needs $(I_2:\pi)=I_1$.
Modulo $x$, the ring is a localization or a power-series
ring over $A/(\pi^3)$, and the preimage of $(\pi^2)$
under multiplication by $\pi$ is exactly $(\pi)$.
Hence $\pi a\in(x,\pi^2)$ implies $a\in(x,\pi)$,
proving the colon equality and exactness.

Let $\delta=\Delta_g(I_1)$, the difference in
(C.1.MF4).
Applying the additive map of the lemma to (C.1.MF5)
gives $\Delta_g(I_2)=-\delta$, while (C.1.MF6)
gives $\Delta_g(I_2)=2\delta$.
Therefore
\[
                              3\delta=0.              \tag{C.1.MF7}
\]
No classification of all maximal Cohen--Macaulay modules
or of the entire singularity Grothendieck group is needed.

There is also an explicit relation of order two.
Let $D=\left(\begin{smallmatrix}0&\Phi\\\Psi&0\end{smallmatrix}\right)$,
$\sigma=\operatorname{diag}(I_2,-I_2)$, and let $D'$
be the following formal polynomial derivative:
its off-diagonal blocks are
\[
 \Phi'=\begin{pmatrix}0&-2\pi\\1&0\end{pmatrix},
 \qquad
 \Psi'=\begin{pmatrix}0&2\pi\\-1&0\end{pmatrix}.
\]
The identities
\[
                  DD'+D'D=3\pi^2I_4,
                  \qquad (D')^2=2\pi I_4             \tag{C.1.MF8}
\]
are polynomial identities over $\mathbb Z[\pi,x,z]$.
They are verified there before specializing the integer
$2$ to $v\pi^e$ in $A$.
In particular no derivation with respect to $\pi$ on
the ramified DVR $A$ is being assumed.

Since $e\geq2$, the element
$\alpha=2/(3\pi^2)=v\pi^{e-2}/3$ belongs to $A$.
Put $h=\alpha D'$ and $F=\sigma-h$.
All of $D,h$ are odd for the displayed grading, so
\[
 TF+FT=2I_4-(Dh+hD)=0,
 \qquad
 F^2=(1+2\pi\alpha^2)I_4.                         \tag{C.1.MF9}
\]
The scalar $1+2\pi\alpha^2$ lies in $1+(\pi)$ and
is a unit.  Thus $F$ is an actual invertible
intertwiner between $T$ and $-T$.

For any free rank-$d$ $C_2$-lattice $M$, tensoring
the admissible augmentation sequence for $R[C_2]$
with $M$ gives
\[
 [M]+[\mathrm{sgn}\otimes M]=d[R[C_2]],
\]
since the diagonal action on $M\otimes_RR[C_2]$
is isomorphic to $d$ regular actions.
Taking $M=(R^4,T)$ and using its sign isomorphism
gives $2\delta=0$.
Together with (C.1.MF7), this proves $\delta=0$.

Tate acyclicity follows from the lemma.
All entries of $\Phi,\Psi$ lie in the maximal
ideal of $R$, so $\bar T=I_4$ on the residue field
$\kappa$ of characteristic two.
If $M$ were projective over $R[C_2]$, its reduction
would be projective, hence free, over the local ring
$\kappa[C_2]=\kappa[t]/(t-1)^2$.
On a nonzero free module over this ring, $t-1$ has
nonzero image, whereas it acts as zero on this
residue representation.  This contradiction proves
the final assertion.
\end{proof}

\noindent\emph{Boundary of an ordinary singularity-class retraction.}
For the construction in (C.1.MF2), the $S_0$-cokernel
of $\widetilde T-I$ fits into an exact sequence
\[
 0\longrightarrow\operatorname{coker}\Psi
  \longrightarrow\operatorname{coker}(\widetilde T-I)
  \longrightarrow\operatorname{coker}\Phi
  \longrightarrow0.
\]
Project to the first coordinate block; injectivity
of $\Phi$ over $S_0$ gives the kernel assertion.
The two endpoint classes sum to zero after quotienting
the maximal Cohen--Macaulay Grothendieck group by free
modules, by the periodic resolution.
Thus the ordinary ungraded class of this Tate cokernel
cannot provide a retraction of (C.1.MF1).
No general failure of injectivity of that map is
asserted from this observation alone.
